\section{(重要)基本知識}
\subsection{高校数学でのベクトル}
高校数学のベクトルと大学数学のベクトルに違いがあるかと聞かれると,本質的には大きな違いはないと私は感じている.
高校数学では「平面ベクトル」「空間ベクトル」までを扱った.
これはつまり,3次元空間までのベクトルということである.それは図で書くことができる最大の次元であるからである.イメージしやすいのがメリットだったが大学数学ではn次元のベクトルまで扱う.しかし,それだけのように感じる.
私が線形代数学を学習するとき,定義はn次元を想定して描かれているが3次元の図を基本イメージして理解に努めていた.4次元ってどうイメージするのか？私にはわからない.

まず高校数学でのベクトルでは$\vec{a}=(1,2,3)$と書いていただろう.ただ,大学では縦書きになる.そして,ベクトルを表していた文字の上にあった矢印はなくなり,代わりに\textbf{太く}書く.\\
つまり,
\[
\vec{a} = (1,2,3) \qquad \rightarrow \qquad \symbfit{a} = \begin{bmatrix}
    1 \\2\\ 3 \end{bmatrix} \text{もしくは} \begin{pmatrix}
        1 \\2\\ 3 \end{pmatrix}\text{となる.}
\]
好みだと思うが,私は[\quad]を使う.注意点は$\begin{vmatrix}1 \\ 2 \\ 3\end{vmatrix}$は間違いですので気をつけるように.これは「行列式」で使う書き方で,意味が変わってしまう.
行列式はSection.\ref{sec:det}で学ぶ.

例えば$s$が実数である場合,実数の集合に属するので$s \in \mathbb{R}$と表現する.
ベクトルも同様に集合の書き方で表現できる.
2次元ベクトル$\symbfit{a}$であれば,$\symbfit{a} \in \mathbb{R}^2$.
3次元ベクトル$\symbfit{a}$であれば,$\symbfit{a} \in \mathbb{R}^3$となる.
\subsection{線形結合}
この言葉は高校数学ではなかった.実はこれは高校生の頃に何気なくやっていることに名称がついただけの話である.
\begin{center}
\underline{線形結合は「k倍したベクトル同士を足せますよ」.これだけの話.}    
\end{center}
\[
c\symbfit{v}+d\symbfit{w}=c\begin{bmatrix}1 \\ 1\end{bmatrix}+d\begin{bmatrix}2 \\ 3\end{bmatrix} =
\begin{bmatrix}c+2d \\ c+3d\end{bmatrix}
\]
この書き方は初学者には特別な変形に見えるかもしれないが,今まであなたがやっていたように
ベクトルを横書きにして( \quad )に書き直して同じ計算をしてみるといい.表記の違いだけということに気づくだろう.
そして上から数字を第1要素,第2要素...なんて読んでいく.これはx軸方向に,y軸方向にとか言っているとz軸方向で詰んでしまうからだと思ってほしい.
そうなると第2要素までのベクトルを2次元ベクトル,第n要素まであるベクトルはn次元ベクトルとなるね.
事前に言っておくが,4次元以降のベクトルは想像できないので,割り切って紙面上数字と思って学習した方が脳に優しいのでおすすめ.
\begin{tcolorbox}[colframe=black!80,colback=white,title=線形結合$_{\text{linear combination}}$]
$c\symbfit{v}$と$d\symbfit{w}$の和は,$\symbfit{v}$と$\symbfit{w}$の線形結合である.
\end{tcolorbox}

必要な道具はここで揃った.実は新しい内容である「\textbf{行列}」に入るにはもうこれで十分である.意外と身構えることではない.
\subsection{行列への拡張}
\paragraph{連立一次方程式をベクトル空間で捉える}
行列の感覚としては最も効率よくベクトルを並べる方法とでも思っている.確かに書き方は変わるが「線形結合」を意識すると難易度はグッと下がる.どんなものか楽しみにしておくとよい.

行列へと自然に持っていくために少し脱線する.皆さんはこの問題は解けるだろうか？
\[
\begin{array}{*{7}{c}}
    x & + &2y & +&3z &= &6 \\
   2x & +&5y & +&2z &= &4 \\
   6x & - &3y & +&z &= &2 \\
\end{array}
\]
おそらく,「3つの式で成り立つx,y,zを求める問題」から「3つの面が交わってできた1点を求める問題」の範囲に収まる認識だと思う.
しかし,線形代数学を学習した学生はそうは考えない.考えてはいけない.なぜならベクトルで考える方がはるかにイメージしやすいからである.
上の式をx,y,zで括ろうと思えばこのように括ることができる.\\
\begin{equation}\label{eq:linear_comb}
x
\begin{bmatrix} 1 \\ 2 \\ 6 \end{bmatrix}
+y
\begin{bmatrix} 2 \\ 5 \\ -3 \end{bmatrix}
+z
\begin{bmatrix} 3 \\ 2 \\ 1 \end{bmatrix}
=
\begin{bmatrix} 6 \\ 4 \\ 2 \end{bmatrix}
\end{equation}
うだうだ書こうとは思ったが,共通因数であるx,y,zを縦に括るともう[ \quad ]で括るほうが見やすく,
そうするともうベクトルの世界に式が入ってしまった.\\
これによって式は,
$\begin{bmatrix} 1 \\ 2 \\ 6 \end{bmatrix}$をx倍,
$\begin{bmatrix} 2 \\ 5 \\ -3 \end{bmatrix}$をy倍,
$\begin{bmatrix} 3 \\ 2 \\ 1 \end{bmatrix}$をz倍したベクトルたちの和が作るベクトルが
$\begin{bmatrix} 6 \\ 4 \\ 2 \end{bmatrix}$となるようにx,y,zを求めたいという式となる.
\begin{figure}[h]
    \centering
    \includegraphics[width=0.6\textwidth]{cal_vec.png}
    \caption{連立方程式のベクトル空間的解釈}
    \label{fig:sample}
\end{figure}\\
※「列ベクトル」,「第k列」はここでは無視で\\
どうだろう？立体空間に傾斜もわからない3つの平面を交わらせて1点をイメージするより,
3本の矢印を伸ばして,その和を右辺のベクトルに持っていく作業と考える方がイメージしやすいとは思わないだろうか？\\
図\ref{fig:sample}にもあるが答えは
$\begin{bmatrix} x \\ y \\ z \end{bmatrix} = 
\begin{bmatrix} 0 \\ 0 \\ 2 \end{bmatrix}$である.
実際に式(\ref{eq:linear_comb})に代入して確かめてみよう.
\[
0
\begin{bmatrix} 1 \\ 2 \\ 6 \end{bmatrix}
+0
\begin{bmatrix} 2 \\ 5 \\ -3 \end{bmatrix}
+2
\begin{bmatrix} 3 \\ 2 \\ 1 \end{bmatrix}
=
\begin{bmatrix} 6 \\ 4 \\ 2 \end{bmatrix}
\]
ん？これ線形結合とは思いませんか？\\
ぬるっと大学数学の「\textbf{線形結合}」と合流する.(そう見えたら狙い通り)
\paragraph{線形結合から行列へ}
ここに来て頑張りどころが初めてくる.ここで行列の書き方と行列の積を紹介する.ここは定義として飲み込んでほしい.
行列同士の計算方法はSection.\ref{sec:matrixRule}で説明するため安心してほしい.
\begin{tcolorbox}[colframe=black!80,colback=white,title=行列とベクトルの積]
\[
\begin{bmatrix}
    & & \\
    \:\symbfit{u}& \symbfit{v} & \symbfit{w} \: \\
    & & \\
\end{bmatrix}
\begin{bmatrix}c \\ d \\ e \end{bmatrix}
=
c\symbfit{u}+d\symbfit{v}+e\symbfit{w},\qquad(\symbfit{u},\symbfit{v},\symbfit{w}\in \mathbb{R}^3,\;c,d,e\in \mathbb{R})
\]
左の[ \quad ]にあるのはベクトルである.今回の問題においては,
\begin{equation}\label{eq:linear_equal}
0
\begin{bmatrix} 1 \\ 2 \\ 6 \end{bmatrix}
+0
\begin{bmatrix} 2 \\ 5 \\ -3 \end{bmatrix}
+2
\begin{bmatrix} 3 \\ 2 \\ 1 \end{bmatrix}
=
\begin{bmatrix} 6 \\ 4 \\ 2 \end{bmatrix}\quad\Longleftrightarrow\quad
\begin{bmatrix}
    1 & 2 & 3 \\
    2 & 5 & 2 \\
    6 & -3 & 1
\end{bmatrix}
\begin{bmatrix}
    0 \\ 0 \\ 2
\end{bmatrix}
=
\begin{bmatrix}
    6 \\ 4 \\ 2
\end{bmatrix}
\end{equation}
\end{tcolorbox}
上の定義とじっくり照らし合わせてほしい.各ベクトルが縦書きを保ったままぎゅっと1つの[ \quad ]にまとまった形をしている.
これを一般的に「\textbf{行列}(matrix)」と呼ばれる.イメージは「行を成す数字の羅列」を略して「行列」であろう(多分).線形結合から行列に書き直すのは慣れればどうってことないが,
初学者にはこの独特な掛け合わせが行列を難しく感じさせる.最初の見慣れた連立方程式から線形結合を経て行列への数字の動きを書いてみるといい.
私も最初は難しく感じた.というのも線形結合を飛ばして行列をみると3\times 3の数字が何を表しているのかすぐにはわからないからだ.
\underline{行列から線形結合.線形結合から行列を繰り返し考えたり書いたりして慣れてみ}\\\underline{よう.}

さて,図\ref{fig:sample}上では挿入されており,「\textbf{列ベクトル}」とある.これについて説明する.Strangの初学本$^{\cite{strang}}$では先に列ベクトルを説明するため先に出ているが,
私個人としては行列を先に見てからの方が納得できると思ったからである.

式(\ref{eq:linear_equal})の行列について話す.行列は横方向に「行」,縦方向に「列」と呼ぶ.
例えば第1行は(1,2,3)であり,第1列は(1,2,6)である.日本人限定の恩恵だが「行」という字の二の部分が方向と一致する.
「列」も同様である.有効な覚え方である.また,(-3)は3行2列の要素なんて言ったりする.改めて図\ref{fig:sample}を見てみよう.これらのベクトルは後に行列の列となる.
そこから「列ベクトル」と名称がつくとなると,なぜ「列ベクトル」なんだという疑問はなくなるのではないだろうか？
今回の3行3列からなる行列を「$3\times3$の行列」と言ったりする.行列が何行何列かを「\hypertarget{size}{\textbf{サイズ}}」という.ちょっとしたことを言うと,列ベクトルも3行1列の行列とも言える.だが,列ベクトル呼びで結構.

3次元ベクトル$\symbfit{a}$であれば,$\symbfit{a} \in \mathbb{R}^3$であった.行列においても集合に属している.
3行2列の行列$A$を挙げるのなら,$A\in \mathbb{R}^{3\times2}$と表現でき,行列$A$は3行2列の実数行列すべての集合の1要素となる.

これから行列主体の内容になっていくが,\textbf{わからなくなったら,左辺の行列と列ベクトルの積は線形結合であることを忘れないでほしい.}
大学の講義における線形代数学のゴールは「対角化(Section.\ref{sec:diagonalization})」であるが,これ本質的に理解させようとすると,この行列の積を線形結合であることに慣れておくことが大事だと考えているからだ.
このテキストでは暗記で終わらせるつもりはない.また,グラムシュミット(Section.\ref{sec:gram})に限らず,これからの行列の計算では「\textbf{わからなくなったら線形結合まで帰着する.}」ことで解決することが多いと信じている.
私自身の現在のレベルまでの学習において,これで乗り越えてきている.
そのため,「列ベクトル」から行列の解説を始めることには\underline{大きな価値があると考えている.}
\newpage
\paragraph{行列に関する用語}
\begin{itemize}
    \item \textbf{n次正方行列：}サイズが$n\times n$の行数と列数が一致した行列を指す.行列は正方形となる.
    \item \textbf{矩形行列：}「くけい」と呼ぶ.正方行列以外の行列.つまり,行数と列数が一致しない行列であり,長方形となる.
\end{itemize}

\textbf{単位行列}は対角線上の要素(対角要素)が1でそれ以外がすべて0となる正方行列を指す.
Strangの一連の本では単位行列は$I$が用いられるが,テストでは$E$が用いられている.例えば,3次単位行列の場合,
\[
    I_3 = \begin{bmatrix}
        1 & 0 & 0 \\ 0 & 1 & 0 \\ 0 & 0 & 1
    \end{bmatrix}
\]

\textbf{零ベクトル・零行列}は
すべての要素がゼロであるベクトルもしくは行列である.
零ベクトルは高校数学のベクトルでは$\vec{0}$であったが,太字で$\mathbf{0}$とする.
また,行列においては単位行列同様慣習的に決まった文字が割り当てられている.行列$O$が用いられる.

\textbf{回転行列}は座標を移動させ,オブジェクトの回転を行うために用いられる.平面オブジェクトであれば2次正方行列となり,立体オブジェクトであれば3次正方行列が用いられる.
回転行列は面白い性質を持つと同時に,これは論理面にも実用面においても重要であるため取り上げる.このテキストでは2次正方行列のみを扱い,行列は$R(\theta)$とする.
\[
    R(\theta) = \begin{bmatrix}
        \cos\theta & -\sin\theta \\ \sin\theta & \cos\theta
    \end{bmatrix}
\]

また,今回は取り上げないが3次の回転行列の場合は$x,y,z$軸回転があるため,3種類紹介する.
3次正方行列と2次回転行列が組み合わさったような非常に合理的な行列である.
\[
    R_x(\theta) = \begin{bmatrix}
        1 & 0 & 0 \\
        0 & \cos\theta & -\sin\theta \\ 
        0 & \sin\theta & \cos\theta
    \end{bmatrix},\qquad
    R_y(\theta) = \begin{bmatrix}
        \cos\theta & 0 & \sin\theta \\ 
        0 & 1 & 0 \\
        -\sin\theta & 0 & \cos\theta 
    \end{bmatrix},\qquad
    R_z(\theta) = \begin{bmatrix}
        \cos\theta & -\sin\theta & 0 \\
        \sin\theta & \cos\theta & 0 \\
        0 & 0 & 1
    \end{bmatrix}
\]

\subsection{演習問題}
\subsubsection*{問題}
\paragraph{問題1.1}
高校数学でのベクトルの問題$^{\cite{FG2}}$を線形代数学の記法に従って解いてみましょう.\\$\vec{p}=(0,1,-1)$を2つのベクトル$\vec{a}=(-2,3,1),\;\vec{b}=(-3,4,2)$と適当な実数$s,t$を用いて,\;$\vec{p}=s\vec{a}+t\vec{b}$の形に表せ.
\begin{enumerate}[label=(\arabic*)]
    \item 3次元ベクトルを[\;]で表現しましょう
    \item 線形結合の式をベクトルを意味するアルファベットで表現しましょう
    \item $\vec{p}$を$s,t$を用いて一つのベクトルで表現しましょう.ただし,ベクトル表記は$[\;]$
    \item 踏まえて$s,t$を求めましょう
\end{enumerate}
\begin{flushright} (自作) \end{flushright}
\paragraph{問題1.2}
次を満たす値$n,k,l$を求めよ.
\begin{enumerate}[label=(\arabic*)]
    \item 5次元ベクトル$\symbfit{r}$において,$\symbfit{r}\in \mathbb{R}^n$を満たす
    \item 4行5列の行列$A$において,$A\in\mathbb{R}^{k\times l}$を満たす.また,この行列を構成する列ベクトル$\symbfit{a}_i$は$\symbfit{a}_i\in \mathbb{R}^n$を満たす.ただし,$1\leq i \leq 5$である.
\end{enumerate}
\begin{flushright} (自作) \end{flushright}
\subsubsection*{解答}
\paragraph{問題1.1}
\begin{enumerate}[label=(\arabic*)]
    \item 高校数学でのベクトルの書き方を紹介した書き方に変換すると,
\[
    \symbfit{p} = \begin{bmatrix}
        0 \\ 1 \\ -1
    \end{bmatrix},\qquad
    \symbfit{a} = \begin{bmatrix}
        -2 \\ 3 \\ 1
    \end{bmatrix},\qquad
    \symbfit{b} = \begin{bmatrix}
        -3 \\ 4 \\ 2
    \end{bmatrix}.
\]

\item $\symbfit{p},\symbfit{a},\symbfit{b}\in \mathbb{R}^3,\;s,t \in \mathbb{R}$とすると,
\[\symbfit{p} = s\symbfit{a} + t\symbfit{b}. \]

\item 
\begin{align*}
    \symbfit{p} &= s\symbfit{a} + t\symbfit{b}\\
    \begin{bmatrix}
        0 \\ 1 \\ -1
    \end{bmatrix}
    &= s\begin{bmatrix}
        -2 \\ 3 \\ 1
    \end{bmatrix}+t
    \begin{bmatrix}
        -3 \\ 4 \\ 2
    \end{bmatrix}\\
    =&
    \begin{bmatrix}
        -2s-3t \\ 3s +4t \\ s+2t
    \end{bmatrix}.
\end{align*}
\item (3)より各行の要素を比較する.
\[
\begin{cases}
    -2s-3t &= 0 \\ 3s+4t &= 1 \\ s+2t &= -1
\end{cases}\qquad s=3,t=-2.
\]
またこれにより,\;$\symbfit{p}$は以下のような線形結合で表現できることがわかる.
\[\symbfit{p} = 3\symbfit{a} -2\symbfit{b} \]
\end{enumerate}
\paragraph{問題1.2}
\begin{enumerate}[label=(\arabic*)]
    \item $n=5$
    \item $k=4,l=5$\\また,\[A = 
    \begin{bmatrix}
        &&&& \\
        &&&& \\
        \symbfit{a}_1&\symbfit{a}_2&\symbfit{a}_3&\symbfit{a}_4&\symbfit{a}_5 \\
        &&&& \\
        &&&& 
    \end{bmatrix}\]
    各列ベクトルは4行なので,$n=4$となる.
\end{enumerate}
\section{線形方程式の求解}
\subsection{(重要)簡約化}
\paragraph{行列方程式}
そういえばベクトルは$\symbfit{a}$と書くなら行列にも何か代表する文字があってもいいはず.あります.小文字のアルファベットが列ベクトルならそれからなる\underline{行列はAとなる.}特に決まりがない場合はAが使われる.\\
先ほどの問題で考える,シンプルに
\begin{equation}
\begin{bmatrix}
    1 & 2 & 3 \\
    2 & 5 & 2 \\
    6 & -3 & 1
\end{bmatrix}
\begin{bmatrix}
    x \\ y \\ z
\end{bmatrix}
=
\begin{bmatrix}
    6 \\ 4 \\ 2
\end{bmatrix}
\qquad \leftrightarrow \qquad A\symbfit{x} = \symbfit{b} \label{eq:Ax_b}
\end{equation}
と書ける.Aは太字じゃないので注意.
\paragraph{簡約化への導入}
さて,私たちはこれから何をするのだろうか？\underline{それは$\symbfit{x}$を求める方法を学ぶことである.}\\
先ほどの問題ではすぐ答えを言っており,解き方の紹介はしていない.答えが予想しやすいので,違和感を感じる人はいなかったかもしれないが,
連立方程式ではこの$\symbfit{x}$にあたる数を求めるのが大変であった.せっかく行列表記にしたのだから,それを活かした新しい画期的な解き方でないとおかしい.

オチを言うと「微妙」であると考えている.私はこれから「\textbf{拡大係数行列}」から始まり,「\textbf{簡約化}」へ触れていくことになるが,計算が辛い.計算ミスもする.私は計算ミスがとても多くこの単元では「良」だった.
計算ミスを減らす方法はあるかと言われると,見つかっていない.しかし,この内容が後の内容につながっていく.具体的には「\textbf{逆行列}」「\textbf{行列式}」「\textbf{余因子行列}」である.
なので,苦行は今後のためのツールだとして頑張って欲しい.
\begin{enumerate}[label=※]
    \item 内情を言うと,行列にすると「\textbf{前進消去と後退代入}」という\underline{画期的な解き方}へアクセスできるが手法自体はテストに出ない.ただ,知っているとテストで早く解けるのでSection.\ref{subsub:Gaussian}(発展)で紹介する.
\end{enumerate}
\bigskip

まず用語から.行列Aは連立方程式の係数からなるので「\textbf{係数行列}」と特に言われる.「拡大係数行列」は字の如く,係数行列を拡大した書き方である.
具体的に行列(\ref{eq:Ax_b})で例を示すと,[A $\symbfit{b}$]こんな感じ.今回の問題では\\
\[
[\,A \:\:\symbfit{b}\,]= 
\begin{bmatrix}
    1 & 2 & 3 & 6 \\
    2 & 5 & 2 & 4 \\
    6 & -3 & 1 & 2
\end{bmatrix},\qquad(A\in \mathbb{R}^{3\times3},\;\symbfit{b}\in\mathbb{R}^3)
\]
つまり行列$A$の要素に$\symbfit{b}$の要素が横にくる行列を拡大係数行列という(「拡大行列」ともいう).なぜそのような書き方をするかであるが,
これから私たちは行列に対して「簡約化」と言うものを行うことで$\symbfit{x}$を求める.拡大係数行列で書いて処理した方が見やすいからである.

「簡約化」は先にゴールを示した方がわかりやすいと思う.簡約化という「処理」を行うことで,どうなるかを示す.
上記にある拡大係数行列を使っても良いのだが,数の並びの都合上見づらいので,新しく拡大係数行列を用意した.簡約化もそれを使って説明する.
\begin{framed}
\begin{equation}    
[\,A \:\:\symbfit{b}\,]=
\begin{bmatrix}
    2& 4 & -2 & 2 \\
    4 & 9 & -3 & 8 \\
    -2 & -3 & 7 & 10
\end{bmatrix}
\quad\rightarrow\quad R=
\begin{bmatrix}
    1& 0 & 0 & -1 \\
    0 & 1 & 0 & 2 \\
    0 & 0 & 1 & 2
\end{bmatrix}\label{eq:ex_gaussian_A}
\end{equation}
\begin{center}
    　　簡約化\\
\end{center}
\vspace{-4mm}
\end{framed}
Aの対角成分が1でそれ以外が0となった.これが簡約化である.ここまでくるとメリットが見えてくる.簡約化された行列は\underline{大変解きやすいのだ.}
例えばこのような連立方程式を解くことを考えてみる.
\[
\begin{array}{ccccccc}
    1x & + &0y & +&0z &= &-1 \\
   0x & +&1y & +&0z &= &2 \\
   0x & + &0y & +&1z &= &2 \\
\end{array}
\]
全ての式が1つの変数に対する1つの値から式となっており,答えはすぐに$x=-1,y=2,z=2$と求まる.\\
この連立方程式を行列で表すと,
\begin{equation}\label{eq:solableByRREF}
\begin{bmatrix}
    1& 0 & 0 \\
    0 & 1 & 0 \\
    0 & 0 & 1 
\end{bmatrix}
\begin{bmatrix}
    x \\ y \\ z
\end{bmatrix}
=
\begin{bmatrix}
    −1 \\ 2 \\ 2
\end{bmatrix}
\end{equation}
さらに拡大係数行列にすると,行列$A$の簡約化後である行列Rと一致する.これによって簡約化によって解きやすくなったことがわかったと思う.

それではその肝心な簡約化の方法であるが,実は皆さんが連立方程式を解くときとそこまで変わらない.中学生でやったようなxとyからなる連立方程式では解き方が2種類あったはず.
1つは$x=,y=$と式変形をして代入する方法.もう1つは「上の式を2倍して,下の式と引く」といった操作を行なって文字を減らす方法.\textbf{簡約化は後者をやっているに過ぎない.}\\

具体的な流れを示す.
簡約化後の拡大係数行列をみると,Aにあたる部分の対角成分が1となっている.この位置の要素をここではpivot(ピボット)と呼ぶ.pivotの定義はこの後出る.対角要素をすべてpivotと呼ぶわけではないので注意.
基本的にはpivotの位置にある要素を1にして,それ以外の要素を0にしていく工程が続く.
\begin{enumerate}[label=(Step.\arabic*), leftmargin=*]
    \item 
まず1行1列は2であるので,第1行を2で割る.
\[
\begin{bmatrix}
    1& 2 & -1 & 1 \\
    4 & 9 & -3 & 8 \\
    -2 & -3 & 7 & 10
\end{bmatrix}
\]
これでpivotは1になった.もう1度簡約化後の拡大係数行列をみると,1行1列の1より下の要素は0なので,0を作る.
第1行を-4倍したものを第2行に足す.(第1行を4倍したものを第2行に引くでも同じ).第1行を2倍したものを第3行に足す.
\[
\begin{bmatrix}
    1& 2 & -1 & 1 \\
    0 & 1 & 1 & 4 \\
    0 & 1 & 5 & 12
\end{bmatrix}
\]
\begin{center}
    \underline{これで第1列の処理は終了}
\end{center}
    \item 
次に第2列の処理をする.2行2列の要素を1にするのだが,今回はたまたまpivotが1になってくれたので上下の要素を0にするだけでよい.
やり方は「式を定数倍して足す」という先程と同じことをする.ポイントは第1列には0があり,いくら式同士でかけたり足したりしても第1列の値は変わらないということ.
つまり,第1列を気にすることなく,第2列の簡約化に集中することができる.
具体的な処理は第2行を2倍して第1行に足す.第2行を−1倍して第3行に足す.
\[
\begin{bmatrix}
    1& 0 & -3 & -7 \\
    0 & 1 & 1 & 4 \\
    0 & 0 & 4 & 8
\end{bmatrix}
\]
    \item 
第3列も同様に処理すると
\[
\begin{bmatrix}
    1& 0 & 0 & -1 \\
    0 & 1 & 0 & 2 \\
    0 & 0 & 1 & 2
\end{bmatrix}
\qquad \text{\textcolor{red}{簡約化完了}}
\]
\end{enumerate}
ここで注意点をいくつか.
\begin{itemize}
    \item 大学の講義では書き方が違うがそちらを推奨する.これは解説用に書き方を変えている.
    \item 今回の簡約化は整数のみとなるが,途中や簡約化後の第4列に分数が現れる方が多い.
    \item 簡約化はRREF行列になるように,\textbf{任意の操作をしても簡約化後の行列は一意になる}.初心者向けに対角成分を1にすることから始めているが,慣れればそうしない方が早い時があることに気づく.
    今回の問題でいうなら,第1行を最初に2で割る前に,第1行を−2倍して第2行へ足す.第1行を1倍して第3行に足す.
    これを先に行った方が素早く0が現れるし,計算ミスが少なくなる.
\end{itemize}
\hspace{1em}
他にも行を入れ替えるなど手段がある.例えば第1行の数字と第2行の数字を交換する.行の入れ替えをしても大丈夫なのか？と思った人もいるかもしれないが,
連立方程式の時を思い出してほしい.式の順番を変えても解には影響ない.連立方程式の形で式を交換する.それを行列にする.これは行の入れ替えをしたという.小技はたくさんある.
つまり,簡約化にはセンスや慣れによって速度や過程の形に差が生まれる.講義や先生が出す解説動画では行列の形で解き方が変わっていくので,
それを見て小技を蓄えていくと良いだろう.
\subsection{(重要)rank,階数なるもの}
rankは簡約化の問題とセットで出題されている.本当は「列空間」(Section.\ref{subsub:colspace}(発展))というものを学んでからの方がrankが何を指しているのかわかりやすい.発展内容であるため,ここでは飛ばして結論だけ言っておしまいにする.
\begin{tcolorbox}[colframe=black!80,colback=white,title=rankとpivotの関係]
\begin{center}
    \textbf{pivotの個数}
\end{center}
\begin{framed}
簡約化後の各行で左から右へ見ると,最初に現れる非零が存在する行がある.その個数.\\\\
注意点としては拡大係数行列でのみrankが存在するわけではない.行列が単体でポンと出題され,「rankを求めよ.」と言われることは十分ある.
簡約化して数えればいい.それで求まる.
\\
今回ならrankは3である.いくつか例を挙げる.
\[
A=
\begin{bmatrix}
    1 & 0 & 0 & 2 \\
    0 & 1 & 0 & -2/5
\end{bmatrix} \qquad
\therefore \;\text{rank}(A) = 2
\]
\[
B=
\begin{bmatrix}
    1 & 0 & 0 & 1 \\
    0 & 1 & 0 & 5 \\
    0 & 0 & 0 & 0
\end{bmatrix} \qquad
\therefore \;\text{rank}(B) = 2
\]
rankの説明で
「簡約化後の各行で左から右へ見ると,最初に現れる非零が存在する行がある.」とあるが,この非零の数をpivot(ピボット)という.非零と言っているが簡約化しているので,ほぼ1のことを言っている.情科の学生なら必ずいつか聞くことになる言葉.
\end{framed}
\end{tcolorbox}
\subsection{(重要)一意に解が定まらない方程式の求解}
前チャプターの「線形方程式の求解」での続きである.これからより広い範囲に対応していく.その問題で求まる解はただ1つに定まっていた.しかし,解が1つに定まらない場合はすでに習っている.
媒介変数とか言っただろうか？ここではパラメータという名前で説明していく.このチャプターでは線形方程式の解き方がより一般的になる.これは1Qの出題範囲でもあり,2Qの出題範囲でもある.
また,これはrankとpivotが深く関わる.忘れた人は復習してから読み進めてほしい.
基本的な計算方法はすでに習得している.拡大係数行列に対して,簡約化を行う.これは解がどのようになろうとも同じである.
そこからどのような解になるかはpivotの位置と個数がそれを決める.つまり解がどのようになるかは簡約化してから知れる物であり,パラメータになろうとも
工程はほとんど変わらないため覚えることもそれほど多くない.

では以下の問題を解いてみよう.\medskip\\
$\begin{cases}
    2x-y+z=4 \\
    x+2y-z=3 \\
    2x-3y+2z=3
\end{cases}$
拡大係数行列を用いて簡約化まで行う.\medskip\\
$[A | \symbfit{b}]=
\begin{bmatrix}
    2 & -1 & 1 & 4 \\
    1 & 2 & -1 & 3 \\
    2 & -3 & 2 & 3
\end{bmatrix}
\rightarrow
\begin{bmatrix}
    1 & 2 & -1 & 3 \\
    0 & -5 & 3 & -2 \\
    0 & 6 & -3 & -3
\end{bmatrix}
\rightarrow
\begin{bmatrix}
    1 & 0 & 0 & 2 \\
    0 & 1 & 0 & 1 \\
    0 & 0 & 1 & 1 
\end{bmatrix}$\\
ここで再び,連立一次方程式に戻すと,
$\begin{cases}
    1x+0y+0z=2 \\
    0x+1y+0z=1 \\
    0x+0y+1z=1
\end{cases}$となり,これはすぐ解がわかる.\\
またもう1つわかることがある.それは解が1つに定まっているということである.\bigskip\\
ではこの問題はどうだろうか？\medskip\\
$\begin{cases}
    x+2y+2z=4 \\
    x+3y+5z=5 \\
    2x+y-5z=5
\end{cases}$\medskip\\
同様に行うと,\\
$[A | \symbfit{b}]=
\begin{bmatrix}
    1 & 2 & 2 & 4 \\
    1 & 3 & 5 & 5 \\
    2 & 1 & -5 & 5
\end{bmatrix}
\rightarrow
\begin{bmatrix}
    1 & 2 & 2 & 4 \\
    0 & 1 & 3 & 1 \\
    0 & -3 & -9 & -3
\end{bmatrix}
\rightarrow
\begin{bmatrix}
    1 & 0 & -4 & 2 \\
    0 & 1 & 3 & 1 \\
    0 & 0 & 0 & 0 \\
\end{bmatrix}$\\
連立一次方程式に戻すと,
$\begin{cases}
    1x+0y-4z=2 \\
    0x+1y+3z=1 \\
    0x+0y+0z=0
\end{cases} \rightarrow
\begin{cases}
    1x-4z=2 \\
    1y+3z=1 \\
\end{cases}
$となる.\\
結論としてはこれはパラメータ(媒介変数)付きの解となる.(試してみるとよい)\medskip\\
この違いはなんだろうか？\\
答えは求める解の個数に対して,\underline{方程式の個数が足りない}ということである.\\
この例では,未知数3個に対してpivotが2個しかないため,1つの変数を自由に動かせる.その結果,解は1つに定まらない.\\
この場合,$x,y,z$のどれかに対してパラメータにすることで解を表現できることになるが,どれをパラメータにすれば良いのだろうか？
$x=1$にして$y$を考える？実は簡約後の拡大係数行列のpivotに注目するとこの問題は比較的容易に解決することができる.
\vspace{1cm}


2つの話をしよう.
\begin{enumerate}
    \item 係数行列$A$と解からなる列ベクトル$\symbfit{x}$を見ると,$A$の各列の要素に変数が掛けられていく結果が方程式となる.解が$x,y,z$とするなら,$A$の第1列には$x$,第2列には$y$という感じに.
    \item $A$全体で見ると,pivotがある列とない列が存在する.行列の簡約化時にはpivotの値が前後の行にある要素を変えることで,簡約化が進んでいる.簡約化の過程では
        pivotが定まってからはpivotは常に変わらず,pivotではない値が変化する.簡約化した$A$は誰が行っても同じ結果になるが過程は異なる.\underline{簡約化は実行前から結果が確定していることが言える}.
\end{enumerate}
つまり,\\
「\textbf{簡約化後に各非零行の最初に現れる1をpivotという.pivot が立つ列に対応する変数を基本変数,pivotが立たない列に対応する変数を自由変数という.}」と言えないだろうか？\\
1つ目の話と私の言いたいことが見えてくる.では第1列にpivot(行は自由)が存在する場合,それに掛けられるxは同様に定数になると考えるのは自然ではないだろうか？\\
列においてpivotは1つだが,連立されているどの方程式も同じ解を示す.どれか1つの方程式でpivotと解が掛けられていれば,別の方程式では掛けられていなくても定数であることは変わらない.
1行1列にpivotがある場合,$x$とpivotが掛けられることになる.$x$の解は定数となる.第2,3行の第1列にはpivotはないが,3つの方程式は同じ解を指すので,xが定数となる.\\
これを踏まえて具体例を見よう.\\

$\begin{cases}
    2x-y+z=4 \\
    x+2y-z=3 \\
    2x-3y+2z=3
\end{cases}$
は簡約化により,
$\begin{cases}
    1x+0y+0z=2 \\
    0x+1y+0z=1 \\
    0x+0y+1z=1
\end{cases}$となった.\\
pivotの位置に注目すると,すべての列に存在する.つまり,すべての解$x,y,z$に対応するpivotが掛けられることを意味する.
確かに解が1つ.つまりすべて定数の解を導き出している.\\

パラメータが出現する問題を見てみよう.\\
$\begin{cases}
    x+2y+2z=4 \\
    x+3y+5z=5 \\
    2x+y-5z=5
\end{cases}$は拡大係数行列の状態では
$
\begin{bmatrix}
    1 & 0 & -4 & 2 \\
    0 & 1 & 3 & 1 \\
    0 & 0 & 0 & 0 \\
\end{bmatrix}$でzが掛けられる列にpivotがない.\\
つまり,$z$が変数(パラメータ)になることがわかる.\\
この問題の最終ゴールを示すなら,$z$がパラメータより$z=t$とする.
それより,$x,y$を$t$を用いて表現すると,$x=4t+2,y=-3t+1$となる.解を線形結合の形で書くと,\\
\[
\begin{bmatrix}
    x \\ y \\ z
\end{bmatrix}=t
\begin{bmatrix}
    4 \\ -3 \\ 1
\end{bmatrix}+
\begin{bmatrix}
    2 \\ 1 \\ 0
\end{bmatrix}
\]\\
これで求められた.\\

\begin{tcolorbox}[colframe=black!80,colback=white,title= パラメータの存在条件]
\begin{center}
   \textbf{行列Aの列数$\neq$pivotの個数}
\end{center}
    \hspace{1cm}$\Longrightarrow\:$\textbf{pivotがない列がパラメータとなる.}\\
    \hspace{1cm}$\Longrightarrow\:$\textbf{pivotがない列数とパラメータの数が一致する.}\\
    \hspace{1cm}$\Longrightarrow\:$\textbf{全ての列にpivotがある場合,解が一意に定まる.(パラメータが存在しない)}
\end{tcolorbox}
特に
正方行列においてすべての列にpivotがある行列を「\textbf{非退化}な行列」という.これはSection.\ref{sub:nondegenerate}で扱うため,読んだ後に振り返るとより理解が深まるかもしれない.
\subsection{(発展)ガウスの消去法と列空間}\label{subsec:Gaussian}
\subsubsection{行列であることのメリット(前進消去と後退代入)}\label{subsub:Gaussian}
テストでは「簡約化せよ」と出題されるので,簡約化されてから,式(\ref{eq:solableByRREF})が示すように線形方程式の解が求められる.加えて,rankも求められた.しかし,簡約化は手間がかかるし,
簡約化しなくても求められたら簡約化したくないものだ.簡約化せずとも求まるのは,
\begin{itemize}
    \item 線形方程式の解$\symbfit{x}$
    \item rank
\end{itemize}
である.つまり,簡約化を強制されてなければ,上記の2つを求める分には\textbf{簡約化は不要なのだ！}
\paragraph{線形方程式の解}
「\textbf{前進消去と後退代入}」の話をしよう.簡約化と手順は大きく変わらないが,「処理」の回数が少なければ計算ミスも減るだろう.
簡約化されたpivotが1でそれ以外が0になった行列をRREF行列(Reduced Row Echelon Form)と言う.名前がカッコいいからと言う理由で紹介する.

この\underline{ガウスの消去法は「前進消去」と「後退代入」の二部構成となっている.}前進消去を行い,後退代入をすることで解が出る.

式(\ref{eq:ex_gaussian_A})で使用した行列を用いる.まず,\textbf{前進消去}
\[
\begin{bmatrix}
    2& 4 & -2 & 2 \\
    4 & 9 & -3 & 8 \\
    -2 & -3 & 7 & 10
\end{bmatrix}
\rightarrow
\begin{bmatrix}
    2& 4 & -2 & 2 \\
    0 & 1 & 1 & 4 \\
    0 & 0 & 4 & 8
\end{bmatrix}
\]
\begin{center}
    　　前進消去
\end{center}
簡約化と同様にゴールを先に示した.簡約化より非零の要素が多く残る形であるが十分である.簡約化との違いは2つ
\begin{enumerate}
    \item \textbf{pivotを1にしなくていい}
    \item \textbf{0にしなくてはいけないのはpivotの下側のみ}
\end{enumerate}
ポイントなのはこの形にするまでの工程は簡約化と同じ工程をするが,pivotを1にする必要がないし,行同士の計算が少ない.注目して欲しいのはpivotが
階段のようにギザギザしている.綺麗な階段にならずとも,行列を斜めに切り,下三角形が全て0になるようにすれば,階段になる.これを「階段化」と言う.
例えば,
\[
\begin{bmatrix}
    1 & 2 & 0 & 3 & 4 \\
    0 & 0 & 1 & 4 & 5 \\
    0 & 0 & 0 & 2 & 1 \\
\end{bmatrix}
\]
これも階段化に成功している.また,これは前進消去が完了した形とも言える.このような行列を「行階段行列」と言います.名前の由来がわかってくれたかな？
\begin{enumerate}[label=※]
    \item「行階段行列」は英語で「REF行列(Row Echelon Form)」という.名前がカッコいいk
\end{enumerate}

また,消去の過程で次操作をしたい行が全て0になってしまったらどうしよう？これはそれより下から都合の良い行を見つけて交換してあげればいい.
都合の良い行とはpivotが存在する行のことである.
「階段」という言葉は同じ階段化は簡約化の行列でも言える.RREF行列は日本語では「簡約階段行列」という.
個人的にはREF行列とRREF行列の違いを細かく知っている必要はないと思うが,説明しておくと,\\
REF行列からpivotを1にし,pivotの上下を0にするとRREF行列となる.

長くなってしまったが\textbf{後退代入}のフェーズを説明する.
初心者向けにまず,前進消去後の行列を連立方程式に戻す(慣れればしなくていい).
\[
\begin{bmatrix}
    2& 4 & -2 & 2 \\
    0 & 1 & 1 & 4 \\
    0 & 0 & 4 & 8
\end{bmatrix}
\rightarrow
\begin{array}{ccccccc}
    2x & + &4y & -&2z &= &2 \\
       &  &1y & +&1z &= &4 \\
       &  &   &  &4z &= &8 \\
\end{array}
\]
見やすいように係数が0になる変数は消している.簡約化と異なり,第1行は変数が3つ生きています.しかし,「後退代入」をすると解ける.
前進消去は上から階段になるように,変数を消していく.恐らく,「前進消去」はここからきているのではないかと思っている.
同様に考えると「後退代入」はREF行列になった行列で下の行列から「\textbf{代入}」していく.つまり第3行から計算する.
第3行では変数が1つなので,$z=2$とすぐ求まる.次に第2行の計算をする.変数が2つだが,第3行で求めたzの値をzに「\textbf{代入}」する.
$y=2$と求まる.最後に第1行のx,yに「\textbf{代入}」すると,\;$x=-1$と求まった.簡約化で求めた解と一致することを確認してみよう.

このように\underline{簡約化して解くより,はるかに少ないステップで解くことが可能となる}.テストでも積極的に活用すれば,簡約化を使用する者より有利に解くことができると考えている.
\begin{enumerate}[label=※]
    \item この一連の処理を「ガウスの消去法(Gaussian Elimination)もしくは「掃き出し法(Row Reduction)」という.
    \item「前進消去(Forward Elimination)」「後退代入(Backward Substitution)」
\end{enumerate}
\paragraph{rankを求める}
\underline{rankはpivotの個数なのでREF行列の状態でも求められる}. 以上\\
つまり,「簡約化」することが条件にされてなければ無理に簡約化する必要はないのである.

例を挙げよう.
\[
\begin{bmatrix}
    1& 2 & -1 & 1 \\
    4 & 9 & -3 & 8 \\
    -2 & -3 & 7 & 10
\end{bmatrix}
\]
こちらは先ほど,簡約化の手順の例でしようした行列である.簡約の手法と同様第1列から順番に処理していく.
\[
\begin{bmatrix}
    1& 2 & -1 & 1 \\
    0 & 1 & 1 & 4 \\
    0 & 1 & 5 & 12
\end{bmatrix}
\]
処理を終えた.簡約化では本来,2行2列の1の上下の要素を0にする.しかし階段化でに則り,第3行のみを0にする.
\[
\begin{bmatrix}
    1& 2 & -1 & 1 \\
    0 & 1 & 1 & 4 \\
    0 & 0 & 4 & 8
\end{bmatrix}
\]
ここで実はrankは3と求まってしまう.

前進消去は用いているpivotより前にある列要素を変えずに,左から右へ操作をしている.前進消去後の行列に注目して欲しい.例えば2行目の行ベクトルを2倍して1行目に対して引く.第3行を1/4倍して...などしていくと,出来る行列はどうなるだろうか?
簡約化が完了する.簡約化はどのような過程を辿っても,簡約化後の行列は同じになる.前進消去を終えたREF行列を見てわかる通り,\underline{RREF行列とpivotの位置は変わっていない}.これが前進消去で十分にrankが求められる理由である.
したがって,\textbf{$\textbf{rank}$を求めるのであれば簡約化までの処理は不要なのである!}
\subsubsection{LU分解}\label{subsub:LU_and_P}
ガウスの消去法はコンピュータ上での求解としても重宝されており,「より優れたアルゴリズムを」と考案されてきた.「数値解析」という授業で触れる予定である.それは\underline{行列$A$を行列$L$と行列$U$の積の形にすること}である.行列同士の積はまだ説明していないが,行列式のpivot公式(Section.\ref{subsub:How_solve_det})で使用するため,紹介する.
また,アルゴリズムやそれを用いた求解方法についてはSection.\ref{sec:LU}で触れるがここでは紹介しない.「数値解析」では授業時間の関係などで内容に限界があるが,テキストは無限に書ける.このアルゴリズムのより詳しい解説に努めるつもりである.そこではTrefethen \& Bau$^{\cite{Trefethen-Bau}}$を参考にしている.

\paragraph{行列Lと行列U}
LU分解(LU decomposition)は行列$L$と行列$U$に分ける.それぞれ,
Lower triangular matとUpper triangularの頭文字から来ており,下三角行列と上三角行列と日本語では訳されている.この名称は行列の要素の配置から来ている.
このLU分解は前進消去を行うために使用した数字が行列$L$に格納され,前進消去後が行列$U$となる.再び式(\ref{eq:ex_gaussian_A})での行列を用いる.
\begin{align*}
    A &= LU \\
    \begin{bmatrix}
        2 & 4 & -2 \\ 4 & 9 & -3 \\ -2 & -3 & 7
    \end{bmatrix} 
    &=
    \begin{bmatrix}
        1 & 0 & 0 \\ 2 & 1 & 0 \\ -1 & 1 & 1
    \end{bmatrix}
    \begin{bmatrix}
    2& 4 & -2 \\ 0 & 1 & 1 \\ 0 & 0 & 4
    \end{bmatrix} \\
    &=
    \begin{bmatrix}
        1 &  &  \\ 2 & 1 &  \\ -1 & 1 & 1
    \end{bmatrix}
    \begin{bmatrix}
    2& 4 & -2 \\  & 1 & 1 \\  &  & 4
    \end{bmatrix}    
\end{align*}
このように要素が上三角形と下三角形からなるように見えるからである.LU分解では三角形の上下は0となるため,3行目のように書かないことが多い.
LU分解の利点は\underline{優れた行列に対してはデータを半分で格納できる}点である.例えば4次正方行列なら,$4^2$個の数字を記憶しておく必要がある.
数千次元の行列レベルになると,10％でもデータを削れればかなりの容量を減らせる.
\paragraph{パスカル行列に対するLU分解}
パスカルの三角形を要素とした行列を「パスカル行列」という.ここでは4次のパスカル行列$\mathrm{pascal}(4)$を考える.
\[
    \mathrm{pascal}(4)=\begin{bmatrix}
        1 & 1 & 1 & 1 \\ 1 & 2 & 3 & 4 \\ 1 & 3 & 6 & 10 \\ 1 & 4 & 10 & 20
    \end{bmatrix}=
    \begin{bmatrix}
        1 & 0 & 0 & 0 \\ 1 & 1 & 0 & 0 \\ 1 & 2 & 1 & 0 \\ 1 & 3 & 3 & 1
    \end{bmatrix}
    \begin{bmatrix}
        1 & 1 & 1 & 1 \\ 0 & 1 & 2 & 3 \\ 0 & 0 & 1 & 3 \\ 0 & 0 & 0 & 1
    \end{bmatrix}
    \qquad(=LU = LL^T)
\]
行列の右上にある$T$は「転置」というものでSection.\ref{subsec:transpose}参照である.
パスカル行列は事前にLU分解して行列$L$を求めておく.\underline{行列$L$のみ覚えておく}.パスカル行列を使用したいときは行列$L$とその転置の積を用いれば再現できる.
つまり,\underline{データが半分となる行列$L$のみで済む}のだ.
\paragraph{対称行列に対するLU分解}
「\textbf{対称行列}」とは対角要素を対称の軸として,要素が対称に並んだ行列を指す.対称行列はかなり重要な行列であり,(2b)で大々的に扱う行列である.また,機械学習や統計,パターン認識でも対称行列は山ほど出現する.
パスカル行列も対称行列である.この対称行列に対するLD分解はより優れている.実例を見せよう.実は何度も使っている式(\ref{eq:ex_gaussian_A})は対称行列である.
その結果は3つの行列の積となる.\newpage
\begin{align*}
    A &= LU \\
    \begin{bmatrix}
        2 & 4 & -2 \\ 4 & 9 & -3 \\ -2 & -3 & 7
    \end{bmatrix} 
    &=
    \begin{bmatrix}
        1 &  &  \\ 2 & 1 &  \\ -1 & 1 & 1
    \end{bmatrix}
    \begin{bmatrix}
        2& 4 & -2 \\  & 1 & 1 \\  &  & 4 \\
    \end{bmatrix} \\
    &=
    \begin{bmatrix}
        1 &  &  \\ 2 & 1 &  \\ -1 & 1 & 1
    \end{bmatrix}
    \begin{bmatrix}
        2 & & \\ & 1 & \\ && 4
    \end{bmatrix}
    \begin{bmatrix}
    1& 2 & -1 \\  & 1 & 1 \\  &  & 1 
    \end{bmatrix}\\
    &= LDL^T
\end{align*}
行列の数がLU分解に対して増えているため,情報が増えたように見えるが,\underline{実質的には変わっていない}.$LDL^T$にすると,\textbf{行列$L$の対角要素は必ず1となる}
ならパソコンの記憶内に覚えておく必要はない.このスペースを有効活用しよう.
それは行列$D$の要素を格納しておくことだ.\textbf{行列$D$は対角要素以外が必ず0となる}.(0は略して記載しない)これは\underline{結果としてLU分解における行列$D$となる}.これが個人的には面白い.行列の積はこの先習うが,n次正方行列同士の積であればの場合は$2n^3-n^2$回四則演算する.つまり計算が大変なのだが,要素をスライドする感覚で,LU分解における行列$U$から$LDL^T$における行列$D$と行列$L$ができてしまう.
\begin{figure}[H]\centering
    \[
    \begin{bmatrix}
        2& 4 & -2 \\  & 1 & 1 \\  &  & 4 \\
    \end{bmatrix}
    \quad \longrightarrow \quad
        \begin{bmatrix}
        1 &  &  \\ 2 & 1 &  \\ -1 & 1 & 1
    \end{bmatrix}
    \begin{bmatrix}
        2 & & \\ & 1 & \\ && 4
    \end{bmatrix}
    \begin{bmatrix}
        1& 2 & -1 \\  & 1 & 1 \\  &  & 1 
    \end{bmatrix}=A
    \]
    \caption{行列$L$さえ覚えておけばすぐ後退代入に使え,行列$A$は再現可能}
\end{figure}
対称行列のような都合の良い行列に関しては半分の記憶容量で済むことがわかった.ちなみにLU分解をするのに必要な計算回数はだいたい$\frac{2}{3}n^3$だが,$LDL^T$ならその半分で済む.
行列演算で対称行列が出現することがどれだけありがたいか垣間見れる一例である.
\subsubsection{列空間とrankの関係}\label{subsub:colspace}
カリキュラムが変わらなければ,線形代数学は「空間」というものに入っていく.これは空間のウォーミングアップになることを願う.
私が学習した内容で$A\symbfit{x}=\symbfit{b}$の考え方を変えてくれた素晴らしい内容である.

線形結合から行列にシフトしたタイミングを思い出して欲しい.$A\symbfit{x}$はAを構成する列ベクトルをそれぞれx倍y倍...としたベクトルたちをイメージする.
そのベクトルの合成によって指すベクトルが$\symbfit{b}$と一致することを式で表すと,$A\symbfit{x}=\symbfit{b}$となる.
では$A\symbfit{x}$が表現できるベクトルに限界はあるのだろうか？わかりやすい例としては,x倍y倍が自然数に限定した場合,明らかに,表現できない領域が現れる.\\
このように,Aの列ベクトルが表現できる領域を「列空間」という.
\begin{tcolorbox}[colframe=black!80,colback=white,title={$A\symbfit{x}=\symbfit{b}$における解の存在条件}]
$A\symbfit{x}=\symbfit{b}$が解を持つのは,$\symbfit{b}$が列空間に含まれるときであり,その時に限る.
\end{tcolorbox}
行列$A$の列空間はC(A)と書く.Cは列の英語の「Column」かと.
では列空間において基本的にx,yは実数全てを取るので,表現できない領域なんてあるのだろうか?と感じる.
列空間は名前からすると,3次元以上の一般的に「空間」と言われるものをイメージすると思うが,Aが平面を作る場合でもそれは列空間となる.

具体的には
\[
C(A)=
x
\begin{bmatrix}1 \\2 \\3\end{bmatrix}
+y
\begin{bmatrix}2 \\ 1 \\ 1\end{bmatrix}
\]
は\underline{3次元の空間に2次元の平面を作る}.この平面も列空間である.実際にやってみて欲しいが,自分がいる部屋で,紙を手で宙に浮かせて欲しい.紙は無限に広いと想像して.$A\symbfit{x}=\symbfit{b}$を満たす場合,
この無限に広がる限られた薄い平面上に基点からの$\symbfit{b}$が存在するということになる.なかなかシビアではないだろうか？※Aが直線を作る場合でも列空間という.

このように$A\symbfit{x}=\symbfit{b}$を満たすのは意外と大変である.そう考えると試験で出題される列空間はなかなか"作られている".現実では収まらないのことが多いらしく,
$\symbfit{b}$が列空間上に存在しないので,列空間上に存在する$\symbfit{b}$に一番近いベクトルを\textbf{近似解}として用いることがある.一番有名なのは\textbf{回帰直線}であるが,この物語の始まりはSection.\ref{sec:startGram}であるのでお楽しみに.
列空間の紹介はこの辺にしておく.\\

本来したかったrankの話をしよう.列空間はベクトルによって作られる空間であるが,その構成するベクトルはなるべく少ない方が情報量が少なくて嬉しい.
例えばこのような列空間はどうだろう?
\[
C(A)=
x
\begin{bmatrix}1 \\2 \\3\end{bmatrix}
+y
\begin{bmatrix}2 \\ 4 \\ 6\end{bmatrix}
+z
\begin{bmatrix}2 \\ 1 \\ 1\end{bmatrix}
\]

この列空間は3つの方向にベクトルが広がって3次元空間を作るように見えないだろうか？答えは間違いであり,正解は平面を作る.
というのもy倍されている列ベクトルはx倍されている列ベクトルの2倍であり,実は同じ方向を指している.つまり,\underline{3つの矢印の内,2つが重複しているので平面となる}.
この考え方でいくと,列空間を表す列ベクトルは無限に増やせてしまい冗長である.言い換えれば,C(A)に必要なのは
$\begin{bmatrix}1 \\ 2\\ 3\end{bmatrix}$と$\begin{bmatrix}2 \\ 1\\ 1\end{bmatrix}$だけで十分であり,この2つでC(A)を十分表現できる.
ではrankと何の関わりがあるのか？実はrankは本来その行列が作る空間を表現する時に\textbf{最低何個の列ベクトルが必要か}を示してくれる.なので,この場合行列Aのrankは2となる.実際に調べてみても良い.「そうなんだ」「確かになる」と思っても,「何故」は残るだろう.
申し訳ないが今はこの程度の情報しか伝えることができない.しかし,これはかなり重要な情報であり,「独立と従属(Section.\ref{sec:depends})」という内容に入った時は必ず学び,ここに戻ってくる.
最終的にはpivotとrankと「独立と従属」がSection.\ref{sub:space_base}で完全に繋がるだろう.小さな豆知識程度のものだとは思わないでほしい.
\subsection{演習問題}
\subsubsection*{問題}
\paragraph{問題2.1}
次のrankを求めよ.\\
\begin{tabular}{llll}
    (1) rank
    $\begin{pmatrix}
        2 & 3 & 1 & 7 \\
        1 & 1 & 1 & 2 \\
        3 & 2 & 4 & 5
    \end{pmatrix}$ 
&(2) rank
    $\begin{pmatrix}
        2 & 1 & 1 & 3 \\
        3 & 2 & 1 & 2 \\
        1 & 1 & 2 & 1 \\
        2 & 1 & 3 & a
    \end{pmatrix}$
&(3) rank$\;E_5$
&(4)$_{\text{難}}$\;rank$\;R(\theta)$
\end{tabular}
\begin{flushright}
    （2025線形代数学1a）改
\end{flushright}
\paragraph{問題2.2}
階数が1となるように,以下の行列の空いているところを埋めよ：
\[
A=
\begin{bmatrix}
    1 & 2 & 4 \\
    2 & & \\
    4 & & 
\end{bmatrix} \qquad
B=
\begin{bmatrix}
    & 9 &  \\
    1 & & \\
    2 & 6 & -3
\end{bmatrix}
\]
\begin{flushright}
    （「ストラング：線形代数イントロダクション」p161）
\end{flushright}
\paragraph{問題2.3}
以下の空欄を適切な数式または数値で埋めよ.パラメーターを要する場合は,文字s,t,u,...をこの順に使用し,「標準的なパラメーターの置き方」に従うこと.\\
(1)連立一次方程式
$\begin{cases*}
    4x+8y=12\\
    7x+12y=21
\end{cases*}$の解は$x=\;$\fbox{\makebox[1.0cm][c]{1}}\;,$y=\;$\fbox{\makebox[1.0cm][c]{2}}\;である.\\
(2)連立一次方程式
$\begin{cases*}
    3x_1+4x_2+x_3+10x_4=15 \\
    2x_1+3x_2+7x_4=11
\end{cases*}$の解は$x_k=\;$\fbox{\makebox[1.0cm][c]{$k$}}\;(k=1,2,3,4)である.
\begin{flushright}
    （2025線形代数学1a）
\end{flushright}
\paragraph{問題2.4}
同次形の連立1次方程式を解け.$\symbfit{x}$を線形結合で表現せよ.\\
\[\begin{pmatrix}
    1 & -2 & 0 & 0 & 1 & 2 \\
    1 & -1 & 1 & 0 & 1 & 0 \\
    2 & -5 & -1 & 0 & 3 & 3 
\end{pmatrix}
\begin{pmatrix}
    x_1 \\ x_2 \\ x_3 \\ x_4 \\ x_5 \\ x_6 
\end{pmatrix}=
\begin{pmatrix}
    0 \\ 0 \\ 0
\end{pmatrix}\]
\begin{flushright}
    （三宅敏恒「線形代数概論」p62）
\end{flushright}
\paragraph{問題2.5}
連立一次方程式
\[\begin{cases*}
    kx+(k^2-3k)z=3k-2 \\
    -2x-y+4z=-7 \hspace*{1.3cm}(k\text{は定数})\\
    x+y-3z=5
\end{cases*}\]
について,以下の空欄を適切な数式または数値で埋めよ.パラメーターを要する場合は,文字s,t,u,...をこの順に使用し,「標準的なパラメーターの置き方」に従うこと.
\begin{itemize}
    \item 解が存在しないならば$k=\;$\fbox{\makebox[1.0cm][c]{1}}\;である.
    \item 解が無数にあるならば$k=\;$\fbox{\makebox[1.0cm][c]{2}}\;で,このとき$x=\;$\fbox{\makebox[1.0cm][c]{3}}\;である.
    \item 解がただ一つならば$z=\;$\fbox{\makebox[1.0cm][c]{4}}\;である.
\end{itemize}
\begin{flushright}
    （2025線形代数学1a）
\end{flushright}
\paragraph{問題2.6}
連立一次方程式が解を持つように,$a,b$を定めよ.\\
\[\begin{pmatrix}
    2 & 1 & 3 \\
    0 & -1 & 1 \\
    1 & 1 & 1
\end{pmatrix}
\begin{pmatrix}
    x_1 \\ x_2 \\ x_3
\end{pmatrix}=
\begin{pmatrix}
    1 \\ a \\ b
\end{pmatrix}\]
\begin{flushright}
    （三宅敏恒「線形代数概論」p64）
\end{flushright}
\subsubsection*{解答}
\paragraph{問題2.1}
\begin{enumerate}[label=(\arabic*)]
    \item 
rank
$\begin{pmatrix}
    2 & 3 & 1 & 7 \\
    1 & 1 & 1 & 2 \\
    3 & 2 & 4 & 5
\end{pmatrix}\rightarrow$
rank
$\begin{pmatrix}
    0 & 1 & -1 & 3 \\
    1 & 1 & 1 & 2 \\
    0 & -1 & 1 & -1
\end{pmatrix}\rightarrow$rank
$\begin{pmatrix}
    1 & 0 & 2 & 1 \\
    0 & 1 & -1 & 1 \\
    0 & 0 & 0 & 2
\end{pmatrix}=3$
\item 
rank
$\begin{pmatrix}
    2 & 1 & 1 & 3 \\
    3 & 2 & 1 & 2 \\
    1 & 1 & 2 & 1 \\
    2 & 1 & 3 & a
\end{pmatrix}\rightarrow$rank
$\begin{pmatrix}
    2 & 1 & 1 & 3 \\
    1 & 1 & 0 & -1 \\
    1 & 1 & 2 & 1 \\
    0 & 0 & 2 & a-3
\end{pmatrix}\rightarrow$rank
$\begin{pmatrix}
    0 & -1 & 1 & 5 \\
    1 & 1 & 0 & -1 \\
    0 & 0 & 2 & 2 \\
    0 & 0 & 2 & a-3
\end{pmatrix}$\\ 
\begin{flushright}
\rightarrow rank $\begin{pmatrix}
    1 & 1 & 0 & -1 \\
    0 & 1 & -1 & -5 \\
    0 & 0 & 1 & 1 \\
    0 & 0 & 0 & a-5
\end{pmatrix}
\quad \therefore
\begin{cases*}
    3 \quad (a=5) \\
    4 \quad (a \neq 5)
\end{cases*}$
\end{flushright}
\item 単位行列はRREF行列であることがわかる.つまり5次の単位行列であるから,rankは5となる.
\item $\cos\theta\neq 0$のとき,2行目から$\dfrac{\sin\theta}{\cos\theta}\times$1行目を引く.\medskip\\
rank $\begin{bmatrix}
    \cos\theta & -\sin\theta \\ \sin\theta & \cos\theta
\end{bmatrix}$\rightarrow rank 
$\begin{bmatrix}
    \cos\theta & -\sin\theta \\ 0 & \cos\theta +\frac{\sin^2\theta}{\cos\theta}
\end{bmatrix}=$ rank 
$\begin{bmatrix}
    \cos\theta & -\sin\theta \\ 0 & \frac{1}{\cos\theta}
\end{bmatrix}$\medskip\\
rankを調べるのであればREF行列までで十分である.よってrankは2である.
また,$\cos\theta=0$のとき,\medskip\\
rank $\begin{bmatrix}
    \cos\theta & -\sin\theta \\ \sin\theta & \cos\theta
\end{bmatrix}=\begin{bmatrix}
    \cos\theta & -\sqrt{1-\cos^2\theta} \\ \sqrt{1-\cos^2\theta} & \cos\theta
\end{bmatrix}=$ rank 
$\begin{bmatrix}
    0 & -1\\ 1 & 0
\end{bmatrix}$\medskip\\
いずれもrankは2である.
\end{enumerate}
\paragraph{問題2.2}
行列のサイズによらず,第1列にのみpivotがあれば良い.\\
Aは第2行,第3行が[1 2 4]の倍数になれば消える.\\
よって
\[A=
\begin{bmatrix}
    1 & 2 & 4 \\
    2 & 4 & 8 \\
    4 & 8 & 16
\end{bmatrix}\]
第2列の1を活かしたい.\\
\[B=
\begin{bmatrix}
    & 9 & \\
    1 & 3 & -3/2 \\
    2 & 6 & -3    
\end{bmatrix}\rightarrow
\begin{bmatrix}
    3 & 9 & -9/2 \\
    1 & 3 & -3/2 \\
    2 & 6 & -3  
\end{bmatrix}\]
\paragraph{問題2.3}
\begin{enumerate}[label=(\arabic*)]
    \item 拡大係数行列より\\
$[A|\symbfit{b}]=
\begin{bmatrix}
    4 & 8 & 12 \\
    7 & 12 & 21
\end{bmatrix}\rightarrow
\begin{bmatrix}
    1 & 2 & 3 \\
    1 & 2 & 3 \\
\end{bmatrix}\rightarrow
\begin{bmatrix}
    1 & 2 & 3 \\
    0 & 0 & 0
\end{bmatrix} \hspace{1cm}\therefore
\begin{cases}
    x+2s=3 \\
    y=s 
\end{cases}\quad \rightarrow \quad
\begin{cases}
    x = -2s+3\\
    y = s
\end{cases}$
\item 拡大係数行列より\\
$[A|\symbfit{b}]=
\begin{bmatrix}
    3 & 4 & 1 & 10 & 15 \\
    2 & 3 & 0 & 7 & 11
\end{bmatrix} \rightarrow
\begin{bmatrix}
    1 & 1 & 1 & 3 & 4 \\
    2 & 3 & 0 & 7 & 11
\end{bmatrix} \rightarrow
\begin{bmatrix}
    1 & 1 & 1 & 3 & 4 \\
    0 & 1 & -2 & 1 & 3 
\end{bmatrix}\rightarrow
\begin{bmatrix}
    1 & 0 & 3 & 2 & 1 \\
    0 & 1 & -2 & 1 & 3 
\end{bmatrix}$\\\\
\hspace{2cm}
$\therefore \begin{cases}
    x_1 +3s+2t=1 \\
    x_2 -2s+t=1
\end{cases}\quad \rightarrow \quad
\begin{cases}
    x_1 = -3s-2t+1 \\
    x_2 = 2s-2t+3 \\
    x_3 = s \\
    x_4 = t
\end{cases}$
\end{enumerate}
\fbox{コメント}

ガウスの消去法を使うことを考えればREF行列までの変形で十分である.しかし,パラメーター付きの解は長く,代入する方が,計算がめんどくさくなる.
この程度の行列ならRREF行列まで変形させてしまうのも手である.\\

\paragraph{問題2.4}
拡大係数行列より\\
$[A|\symbfit{b}]=
\begin{bmatrix}
    1 & -2 & 0 & 0 & 1 & 2 \\
    1 & -1 & 1 & 0 & 1 & 0 \\
    2 & -5 & -1 & 0 & 3 & 3 
\end{bmatrix} \rightarrow
\begin{bmatrix}
    1 & -2 & 0 & 0 & 1 & 2 \\
    0 & 1 & 1 & 0 & 0 & -2 \\
    0 & -1 & -1 & 0 & 1 & -1 
\end{bmatrix}\rightarrow
\begin{bmatrix}
    1 & -2 & 0 & 0 & 1 & 2 \\
    0 & 1 & 1 & 0 & 0 & -2 \\
    0 & 0 & 0 & 0 & -1 & -3
\end{bmatrix}\\ \hspace{1cm}\rightarrow
\begin{bmatrix}
    1 & -2 & 0 & 0 & 0 & 1 \\
    0 & 1 & 1 & 0 & 0 & -2 \\
    0 & 0 & 0 & 0 & -1 & -3    
\end{bmatrix}\quad \therefore
\begin{cases*}
    x_1+2x_2+x_6=0\\
    x_2+x_3-2x_6=0 \\
    x_5-3x_6=0
\end{cases*}$\\
\qquad$\symbfit{x}=c_1
\begin{pmatrix}
    -2 \\ -1 \\ 1 \\ 0 \\ 0 \\ 0
\end{pmatrix}+c_2
\begin{pmatrix}
    0 \\ 0 \\ 0 \\ 1 \\ 0 \\ 0
\end{pmatrix}+c_3
\begin{pmatrix}
    -1 \\ 2 \\ 0 \\ 0 \\ 3 \\ 1
\end{pmatrix} \quad (c_1,c_2,c_3 \in \mathbb{R})$\\
\paragraph{問題2.5}
拡大係数行列より\\
$[A|\symbfit{b}]=
\begin{bmatrix}
    k & 0 & k^2-3k & 3k-2 \\
    -2 & -1 & 4 & -7 \\
    1 & 1& -3 & 5
\end{bmatrix}\rightarrow
\begin{bmatrix}
    k & 0 & k^2-3k & 3k-2 \\
    0 & 1 & -2 & 3 \\
    1 & 1 & -3 & 5
\end{bmatrix}\rightarrow
\begin{bmatrix}
    k & 0 & k^2-3k & 3k-2 \\
    0 & 1 & -2 & 3 \\
    1 & 0 & -1 & 2
\end{bmatrix}\\ \hspace{3cm}\rightarrow
\begin{bmatrix}
    1 & 0 & -1 & 2 \\
    0 & 1 & -2 & 3 \\
    0 & 0 & k(k-2) & k-2
\end{bmatrix}$
\begin{enumerate}[label=(\arabic*)]
    \item 解が存在しない時
$\begin{bmatrix}
    0 & 0 & 0 & \text{実数}(\neq 0)
\end{bmatrix}$となる.つまり$k=0$このとき,
\[
[A|\symbfit{b}] =
\begin{bmatrix}
    1 & 0 & -1 & 2 \\
    0 & 1 & -2 & 3 \\
    0 & 0 & 0 & 2
\end{bmatrix}\]
\item 解が無数にある時rank($[A|\symbfit{b}]) < $列数の時であるから第3行がすべて0となる$k$を考える.つまり$k=2$このとき
\[
[A|\symbfit{b}] =
\begin{bmatrix}
    1 & 0 & -1 & 2 \\
    0 & 1 & -2 & 3 \\
    0 & 0 & 0 & 0
\end{bmatrix}
\]
より$x=s+2$
\item 線形方程式の解は3パターンで解なし,無数解,唯一解である.前2つはk=0,2の時なので,$k \neq 0,2$を考えるのが妥当である.
唯一解の時3行3列が1となればいいので,第3行を$k(k-2)$で割ると,
\[
[A|\symbfit{b}] =
\begin{bmatrix}
    1 & 0 & -1 & 2 \\
    0 & 1 & -2 & 3 \\
    0 & 0 & 1 & \frac{1}{k}
\end{bmatrix}.
\]
後退代入より$\displaystyle z=\frac{1}{k}$
\end{enumerate}
\fbox{コメント}

rankが絡む問題はREF行列を活用すべき.REF行列から解を求める時はもちろん後退代入.\\なぜならREF行列への変形が前進消去なのだから.\\
\paragraph{問題2.6}
拡大係数行列より\\
$[A|\symbfit{b}]=
\begin{bmatrix}
    2 & 1 & 3 & 1\\
    0 & -1 & 1 & a\\
    1 & 1 & 1 & b
\end{bmatrix}\rightarrow
\begin{bmatrix}
    0 & -1 & 1 & 1-2b \\
    0 & -1 & 1 & a \\
    1 & 1 & 1 & b
\end{bmatrix} \rightarrow
\begin{bmatrix}
    1 & 1 & 1 & b \\
    0 & -1 & 1 & a \\
    0 & 0 & 0 & 1-2b-a
\end{bmatrix}$\\
解を持つので$1-2b-a=0.$
\section{(重要)行列操作のルール}\label{sec:matrixRule}
\subsection{和と差とスカラー倍の法則}
$A\symbfit{x}=\symbfit{b}$で行列同士の掛け算を少しやっているが,改めてまとめる.これから行列同士の計算が本格的に必要になってくるからだ.
\begin{tcolorbox}[colframe=black!80,colback=white,title={行列和と行列差の計算可能な条件,スカラー倍の計算}]
和と差について,\textbf{行列の\hyperlink{size}{サイズ}が同じであれば}計算することができる.
\[
\begin{bmatrix}
    1 & 2 \\ 3 & 4 \\ 0 & 0
\end{bmatrix}
+
\begin{bmatrix}
    2 & 2 \\ 4 & 4 \\ 9 & 9
\end{bmatrix}
=
\begin{bmatrix}
    3 & 4 \\ 7 & 8 \\ 9 & 9
\end{bmatrix}
\]
また,行列に対して何倍かする時,その係数は\textbf{すべての要素に施される}.
\[
2\begin{bmatrix}
    1 & 2 \\ -3 & 4 \\ 0 & 5
\end{bmatrix}
=
\begin{bmatrix}
    2 & 4 \\ -6 & 8 \\ 0 & 10
\end{bmatrix}
\]
\end{tcolorbox}
Aを-Aに足すと,すべての要素は0である「\textbf{零行列}」となる.
簡単だと感じただろうか？しかし,本格的な問題は「行列積」.つまり行列同士の掛け算である.
これからそれについて触れる前に要素の示し方について触れておこう.

i行j列にある要素は$a_{ij}$また$A(i,j)$と書くのが一般的である.
例えば,第1行のn個の要素は$a_{11},a_{12}, \cdots ,a_{1n}$となる.行列の左下の要素は$a_{m1}$となり,右下は$a_{mn}$となる.
\subsection{積の法則}
\subsubsection{基本操作}
\paragraph{要素同士の掛け合わせ手順}
では行列積の話をしよう.行列積には和と差同様,行列の形によっては計算できないときがある.
\begin{tcolorbox}[colframe=black!80,colback=white,title=行列積の計算可能なサイズの組み合わせ]
行列$A$,Bが計算できるためには,「Aがn列からなるとき,Bはn行からなる必要がある」\\
文字で見ると難しいが視覚的にわかりやすい書き方がある.
\[A:m\times n \; \; B:n\times k\]
このとき,\underline{隣り合う数字が同じなら計算可能}\\
この場合,nのことである.
また,行列積の結果は,隣り合う数字が消されて結合されたように,$m\times k$となる.
\end{tcolorbox}
例えば,
\[
\begin{bmatrix}
    1 & 2 \\ 3 & 4 \\ 0 & 0
\end{bmatrix}
\begin{bmatrix}
    2 & 2 \\ 4 & 4 \\ 9 & 9
\end{bmatrix}=?\qquad(\text{※行列同士の積は}\times\text{や}\cdot\text{を書かない})
\]
は計算ができない.何故なら,$3\times\underline{2}$と$\underline{3}\times2$で隣り合う数字が異なるからである.
しかし,
\[
\begin{bmatrix}
    1 & 4 \\ 1 & 5 \\ 4 & 2
\end{bmatrix}
\begin{bmatrix}
    3 & 2 \\ 1 & 0
\end{bmatrix}
\]は計算が可能である.サイズを見ると,隣り合う数字が2であるからである.

計算方法であるが,まず授業で紹介される手法は好まない.
私が勝手に思っているかもしれないが,計算結果を書いていく時,掛け合わせのミスは無くなるだろうが,純粋な和の結果を描くときに計算ミスをする可能性があるなと感じている.
和の計算に限っては暗算になるからである.問題点は和が正しいか検算する時に途中式がないので,再度計算し直す必要があるということである.小さい結果になるならそんな話はしないが,分数同士や2桁になってくるとまずい.時間がかかる.
そんなこと気にする人の方が少ないかもしれないし,
気にしない人は無視していいのだが,計算ミスが多い私としては工夫したいところではある.先どり学習で読んでいるのであれば気にしなくて良い.以降に示す手法の使用を一考してみてはどうだろうか？\medskip\\
\fbox{ポイント}
\begin{framed}
「列ベクトル\times 行ベクトル」という方法でも行列積は求まる.\\
\begin{equation}AB=
\begin{bmatrix}
    1 & 4 \\ 1 & 5 
\end{bmatrix}
\begin{bmatrix}
    3 & 2 \\ 1 & 0
\end{bmatrix}
=
\begin{bmatrix}
    1 \\ 1
\end{bmatrix}
\begin{bmatrix}
    3 & 2
\end{bmatrix}
+
\begin{bmatrix}
    4 \\ 5
\end{bmatrix}
\begin{bmatrix}
    1 & 0
\end{bmatrix}
=
\begin{bmatrix}
    3 & 2 \\ 3 & 2
\end{bmatrix}
+
\begin{bmatrix}
    4 & 0 \\ 5 & 0
\end{bmatrix}
=
\begin{bmatrix}
    7 & 2 \\ 8 & 2
\end{bmatrix}\label{eq:ex_col_times_row}
\end{equation}
授業でやった計算の組み合わせとは本質的には変わっていない.
\end{framed}
上記の計算はこのような法則に則っている.
\[
\begin{array}{c|cc}
     & 1 & 0 \\ \hline
  4  & 4\times1 & 4\times0 \\
  5  & 5\times1 & 5\times0 \\
\end{array}
\]
列ベクトルと行ベクトルを表側と表頭におき,各列を対象に各行と掛算を行う.これは行列積が可能なサイズであれば任意のサイズで成り立つ.
計算がシンプルに見え,かつ工程が多く見れるので検算しやすいのではと感じている.これは好みの問題だが,私はこのやり方がやりやすい.

しかし,\textbf{必ず注意して欲しいのは,行\times 列ではない.}つまり,
\[
\begin{bmatrix}
    1 & 4 
\end{bmatrix}\begin{bmatrix}
    3 \\ 1
\end{bmatrix}+\begin{bmatrix}
    1 & 5
\end{bmatrix}\begin{bmatrix}
    2 \\ 0
\end{bmatrix}
\]
は間違いである.ただ,実際にこれは計算すると誤っていることに気づく.\\
\underline{行列のサイズによらず,行\times 列をすると必ず,結果は1\times1の行列となる.}理由はシンプルで
1行m列とm行1列の計算は$1\times1$の行列となるからである.
\paragraph{行列同士の掛け合わせ手順}
上でその行列同士の積が取れるかどうかはサイズによると説明した.そのため,今までやってきたような$3\times4=12$のような掛け算で行えた手法が使えなくなる.このセクションではそれについて説明する.短い内容だが,\textbf{行列同士の計算の根底となるため非常に重要である}.
行列$A,B,C$において,
\[
    AB=CD
\]
が成りたっているとする.ただし,$A$のサイズは$n\times k$とする.\\
このとき,$m\times n$の行列$E$を掛けるとき,サイズ的に$EA$と積はとれる.
\[
    EAB=ECD
\]
と計算可能である.では,$3\times4=12$のように順番を並び替えて,
\[
    EAB=CED,\quad EAB=CDE
\]
などとしていいのだろうか？答えはバツである.理由は単純で,サイズの組み合わせで計算できない\underline{可能性がある}からである.例えば,行列$D$のサイズが$j \times p$なら$p= m$とは限らない.
\[
\text{行列を両辺で掛けるときは,\textbf{左から},\textbf{右から}である.}
\]

\[
AB = C \quad\Longrightarrow\quad
\begin{cases}
EAB = EC\\
ABE = CE
\end{cases}\\
\text{※いずれもサイズ的に計算可能と仮定する}
\]

では,\underline{可能性がある}と書いたが,サイズの組み合わせによって計算できたら順番を並び替えてもいいのか？と思うかもしれない.
実はそうではない.入れ替えしても問題ない行列は$AB=BA$が成り立つ行列に限り可能であり,これを可換という.可換が成り立つ代表例はあとに紹介する逆行列かもしれない.しかし,ここまで学ぶ必要は無いと思うし,可換が成り立つかもわからない行列に対して無理に入れ替えして間違えるのは危険である.
\subsubsection{(発展)単位行列と置換行列の融合}\label{subsub:I_and_P}
個人的には発展といいつつ,\underline{関連する内容の見え方が大きく変わり感動した内容}ため,ぜひ読んで欲しい内容である.
\paragraph{スカラー倍の単位行列}
単位行列はいくら行列に対して掛けても要素を変えない.これは演習問題3.2 (1)で確認できる.
\[
    I^nA = AI^n = A\qquad(\text{ただし,行列積がとれるサイズを満たしている場合})
\]
これは行列積をシグマを用いて表現するとよくわかる.$i$行$j$列の単位行列の要素を$I_{ij}$とすると,
\begin{equation}
    I_{ij} = \begin{cases}
        1 \qquad(i=j)\\
        0 \qquad(i\neq j)
    \end{cases}\label{eq:I}
\end{equation}
同様に行列$A$の要素を$A_{ij}$とし,行列積後の行列$IA$の要素を$(IA)_{ij}$とする.
式(\ref{eq:ex_col_times_row})における1行1列は
\begin{itemize}
    \item 行列$A$の1行1列と行列$B$の1行1列の積
    \item 行列$A$の1行2列と行列$B$の2行1列の積
\end{itemize}
の和から成り立っている.他の要素も調べていくと結果として,
\begin{equation}
    (AB)_{ij} = \sum_k A_{ik}B_{kj} \label{eq:sum_AB}
\end{equation}
であることがわかる.今回の場合,
\[
    (IA)_{ij} = \sum_k I_{ik}A_{kj}
\]
となるが,式(\ref{eq:I})より$k=i$のときしか項が残らない.つまり,
\[
    (IA)_{ij} = A_{ij}
\]
となり,\underline{単位行列は確かに要素を変えない}ことが確認できた.

コレがわかると全ての要素$n$倍する行列$M$というのも作れる.それは言うまでもなく,
\[
    M_{ij} = \begin{cases}
        n \qquad(i=j)\\
        0 \qquad(i\neq j)
    \end{cases}
\]
である.しかし単位行列が要素を変えないことがわかればシグマは不要である.例えば$n=4$の3次正方行列においては
\[
    \begin{bmatrix}
        4 & & \\
         & 4 & \\ 
         & & 4
    \end{bmatrix}
    A = 
    4\begin{bmatrix}
        1 & & \\ & 1 & \\ && 1
    \end{bmatrix}A=
    4IA
    =4A
\]
これにより行列$A$のすべての要素を4倍する行列であることがわかる.

次に一般化しての紹介では難しいので,以下のような3次正方行列を用いて考える.
\[
    D = \mathrm{diag}(2,4,5) = \begin{bmatrix}
        2 & 0 & 0 \\ 0 & 4 & 0 \\ 0 & 0 & 5
    \end{bmatrix}
\]
$\mathrm{diag}$は対角行列を表現する書き方であるが,これは先程の行列$M$における各行で要素数が異なる場合を紹介するために用意した行列である.
再びシグマを使う.式(\ref{eq:sum_AB})における$k=i$の項しか残らない流れは変わらないため,
\[
    (DA)_{ij}=D_{ii}A_{ij},\qquad(D_{11}=2,D_{22}=4,D_{33}=5)
\]
となる.なので,
\[
\begin{aligned}
    (DA)_{1j} &= 2A_{1j}\\
    (DA)_{2j} &= 4A_{2j}\\
    (DA)_{3j} &= 5A_{3j}
\end{aligned}
\quad \longrightarrow \quad 
DA = \begin{bmatrix}
    2A_{1j} \\ 4A_{2j} \\ 5A_{3j}
\end{bmatrix}
\]
※$A_{1j}$は第1行は固定で列番号が変わるので,行列$A$第1行ベクトルという意味となる.\medskip\\
つまり\underline{対角行列は各行に対して要素を施している}ことがわかる.
行列積における要素の組み合わせは比較的複雑だが,対角行列に関してはシンプルで単位行列から進めていくと意外と奥が深いことがわかる.
これらを知っていると行列をより自由に操作出来るようになるだろう.
\paragraph{置換行列}
「スカラー倍の単位行列」において対角行列との積がどのように行列に作用するかを見た.これでも十分有益な情報だが,\textbf{置換行列$P$}の導入にもなる.
\underline{置換行列は施す行列の行ベクトル同士をそのまま\textbf{交換}させ}\\\underline{ることが可能である}.
置換行列は「置換」の知識が必要であり「離散数学」で習うことになるだろう.線形代数概論$^{\cite{linear_math}}$では行列式の導入で紹介され,
置換$\sigma$が
\[
1 \mapsto k_1,\quad
2 \mapsto k_2,\quad
\cdots,\quad
n \mapsto k_n,\quad
\]
という写像のとき,置換$\sigma$を
\[
    \sigma = \begin{pmatrix}
        1 & 2 & \cdots & n \\ k_1 & k_2 & \cdots & k_n
    \end{pmatrix}
\]
と表している.が\underline{ここまで踏み入っていないため,私なりの直感的な理解で紹介しようと思う}.

置換行列$P$は単位行列の行が入れ替わった形である.
\[
P = \begin{bmatrix}
    1 & & \\
    & 1 & \\
    & & 1
    \end{bmatrix},\quad
\begin{bmatrix}
    1 & & \\
    & & 1 \\
    & 1 & 
\end{bmatrix},\quad
\begin{bmatrix}
    & 1 & \\
    1 & & \\
    & & 1 
\end{bmatrix},\quad
\begin{bmatrix}
     & 1 & \\
    &  & 1 \\
    1 & & 
\end{bmatrix},\quad
\begin{bmatrix}
     &  & 1\\
    1 & &  \\
     & 1 & 
\end{bmatrix},\quad
\begin{bmatrix}
     & & 1\\
     & 1 & \\
    1& &  
\end{bmatrix}
\]
これは3次の置換行列の\underline{全て}である.全部で$3!$個ある.これに関しては置換$\sigma$を使うとわかりやすい.
数字が行番号だとしたら,
\[
    \sigma = \underline{\begin{pmatrix}
        1 & 2 & 3\\ 
        1 & 3 & 2
    \end{pmatrix}},
    \begin{pmatrix}
        1 & 2 & 3\\ 
        2 & 1 & 3
    \end{pmatrix},
    \begin{pmatrix}
        1 & 2 & 3\\ 
        2 & 3 & 1
    \end{pmatrix},
    \begin{pmatrix}
        1 & 2 & 3\\ 
        3 & 1 & 2
    \end{pmatrix},
    \begin{pmatrix}
        1 & 2 & 3\\ 
        3 & 2 & 1
    \end{pmatrix}
\]
例えば下線部の1つ目は,
1行目はそのままで2行目と3行目を交換したことを意味する.全部で3行あるので${}_3P_2=3!$というわけである.

一例として3つ目の置換行列$P$を用いる.行列$A$を三次正方行列として式(\ref{eq:sum_AB})より,
\[(PA)_{ij}=\sum_{k=1}^{3}P_{ik}A_{kj}\]
となる.例えば1行目においては,
\[P_{1j}=\begin{bmatrix}0 & 1 & 0\end{bmatrix}\]
であるから,
\begin{align*}
(PA)_{1j} & = \sum_{k=1}^{3}P_{1k}A_{kj} \\
          & = P_{11}A_{1j} + P_{12}A_{2j} + P_{13}A_{3j} \\
          & = 0\cdot A_{1j} + 1\cdot A_{2j} + 0\cdot A_{3j} \\
          & = A_{2j}
\end{align*}
つまり,\underline{$A_{1j}$が$A_{2j}$になった}ということである.これは行交換である.
\[
\begin{bmatrix}
    & 1 & \\
    1 & & \\
    & & 1 
\end{bmatrix}\begin{bmatrix}
    A_{1j}\,\text{(第1行ベクトル)}\\
    A_{2j}\,\text{(第2行ベクトル)}\\
    A_{3j}\,\text{(第3行ベクトル)}
\end{bmatrix}=
\begin{bmatrix}
    A_{2j}\,\text{(第2行ベクトル)}\\
    A_{1j}\,\text{(第1行ベクトル)}\\
    A_{3j}\,\text{(第3行ベクトル)}
\end{bmatrix}
\]
これは\underline{単位行列から置換行列へ行が交換された行と行列$A$の行ベクトルの交換が一致する}.使用した置換行列$P$は単位行列の2行目1行目が交換されたものである.
これによって行列$A$は2行目と1行目が交換される.

置換全体の集合は\textbf{閉じている}.行の交換後にさらに交換しても,3行しかないベクトルが行える交換の候補の範囲からでることは無いため,$3!$通りのいづれかの行の位置になる.
例えば,
1行目と3行目を交換し,
1行目と2行目を交換し,
2行目と3行目を交換した結果は
1行目と2行目を交換した行の位置となる.
\[ABC=
\begin{bmatrix}
    1&& \\ &&1 \\ &1&
\end{bmatrix}
\begin{bmatrix}
    &1& \\ 1&& \\ &&1
\end{bmatrix}
\begin{bmatrix}
    &&1 \\ &1& \\ 1&&
\end{bmatrix}
=
\begin{bmatrix}
    &1& \\ 1&& \\  &&1
\end{bmatrix}
\]
\begin{enumerate}[label=※]
    \item 行列$C$から始め,左から順に行列を掛けていることに注意せよ.
\end{enumerate}
1行目と2行目を交換し,
1行目と2行目を交換すれば元に戻るため行交換をしていないことと等しい.
\begin{equation}
    \begin{bmatrix}
        &1& \\ 1&& \\ &&1
    \end{bmatrix}
    \begin{bmatrix}
        &1& \\ 1&& \\ &&1
    \end{bmatrix}
    = \begin{bmatrix}
        1&& \\ &1& \\ &&1
    \end{bmatrix}=I_3
\end{equation}
\paragraph{改めてLU分解を見る}
簡約化では行変換を行っていたが,前進消去においても必要に応じて行変換を要する場合がある.そこで置換行列を用いる.
\begin{tcolorbox}[colframe=black!75,colback=white,title=行変換を含むLU分解]
\vspace{-5mm}
\begin{equation}
    PA = LU \label{eq:PA_LU}
\end{equation}
\end{tcolorbox}
行の交換は前進消去の過程で必要になるとわかるとそこでするため,事前に行列$A$に対して適切な置換行列$P$を用意しておくのは,未来予知に近いのは理解できる.注目して欲しいのは
行変換は複数回行われる可能性があるのに式(\ref{eq:PA_LU})を見ると置換行列は1回しか掛けられていない.それで十分なのである.なぜなら前述したように\underline{置換全体の集合は閉じている}からである.
複数回の行変換は1回の行変換と一致する.
\subsection{転置(transpose)}\label{subsec:transpose}
転置とは\textbf{行と列の位置関係を交換すること}である.本によって行列に対して「転置」をしたことを示す書き方が異なる.
\begin{enumerate}
    \item ${}^t\!A$\;( \LaTeX では\verb|{}^t\!A|と書く)
    \item $A^t$(小文字)
    \item $A^T$(大文字)
\end{enumerate}
本テキストでは$A^T$と表現している.

例えば,$3\times5$の行列$A$を転置した$A^T$は$5\times3$となる.
\[
A=\begin{bmatrix}
    11 & 12 & 13 & 14 & 15 \\
    21 & 22 & 23 & 24 & 25 \\ 
    31 & 32 & 33 & 34 & 35 
\end{bmatrix}\qquad\to\qquad
A^T=
\begin{bmatrix}
    11 & 21 & 31 \\
    12 & 22 & 32 \\
    13 & 23 & 33 \\
    14 & 24 & 34 \\
    15 & 25 & 35 
\end{bmatrix}
\]
また,列ベクトルに対してももちろん転置は可能である.
\[
    \symbfit{a}=\begin{bmatrix}
        1 \\ 2 \\ 3
    \end{bmatrix}\qquad\to\qquad
    \symbfit{a}^T=\begin{bmatrix}
        1 & 2 & 3
    \end{bmatrix}
\]
\bigskip\\

\underline{転置に対しては一般的に覚えておくべき重要な法則が存在する}ため,紹介する.
\begin{itemize}
    \item 行列同士の足し算の結果を転置するのと,転置してから足し算しても足される数同士の組み合わせは変わらないので,
\[
    (A+B)^T=A^T+B^T
\]
\item 和と同様,積を転置するのと,転置して積をとっても,組み合わせは変わらない.しかし,転置すると行列のサイズが変わり,積が取れなくなることがある.\\
例を挙げる.$3\times4$と$4\times2$の積において転置すると,$4\times3$と$2\times4$でたしかに積が取れない.そこで,場所を入れ替えると$2\times4$と$4\times3$なって積が取れる.
\[
    (AB)^T = B^TA^T
\]
\end{itemize}
実際に簡単な行列を作って検証してみると一層理解が深まるだろう.また,例では2つの行列だったが,n個の行列でも成り立つ.
\begin{tcolorbox}[colframe=black!80,colback=white,title=転置の性質]
\begin{align*}
    (A_1+A_2+\dots+A_n)^T&=A_1^T+A_2^T+\dots+A_n^T\\
        (A_1A_2\dots A_n)^T&=A_n^T\dots A_2^TA_1^T
\end{align*}
\end{tcolorbox}
\subsection{演習問題}
\subsubsection*{問題}
\paragraph{問題3.1}
行列の四則演算を演習しましょう.定義上計算できない場合は「計算不可能」と答えよ.\\
\begin{tabular}{lll}
    (1)\;$3\begin{bmatrix}
        3 & 2 \\ 5 & 1
    \end{bmatrix}-\begin{bmatrix}
        1 & 2 \\ 3 & -1
    \end{bmatrix}$ & 
    (2)\;$\begin{bmatrix}
        1 \\ 5
    \end{bmatrix}-2\begin{bmatrix}
        1 & 2 \\ -4 & 5
    \end{bmatrix}$ &
    (3)\;$\begin{bmatrix}
        2 & 1 & -3 
    \end{bmatrix}-\begin{bmatrix}
        1 & 1 & 1
    \end{bmatrix}$
\end{tabular}
\begin{flushright}
    (自作)
\end{flushright}
\paragraph{問題3.2}
行列同士の積を演習してみましょう.同様に「計算不可能」と答えよ.\\
\begin{tabular}{lll}
    (1)\;$\begin{bmatrix}
        1 & 2 & 3 \\ 4 & 5 & 6
    \end{bmatrix}\begin{bmatrix}
        1 & 0 & 0 \\ 0 & 1 & 0 \\ 0 & 0 & 1
    \end{bmatrix}$ & 
    (2)\;$\begin{bmatrix}
        1 & -1 & 3 \\ 0 & 1 & -6 \\ 7 & 8 & 4
    \end{bmatrix}\begin{bmatrix}
        -3 & 2 & 0 \\ 3 & 1 & 1 
    \end{bmatrix}$ & 
    (3)\;$R(\theta)R(\phi)$
\end{tabular}
\begin{flushright}
    (自作)
\end{flushright}
\paragraph{問題3.3}
\[A = \begin{bmatrix}
    1 & 2 \\ 2 & 4
\end{bmatrix},\qquad
B = \begin{bmatrix}
    4 & 0  \\ -1 & 2 
\end{bmatrix}\]
を用いて,転置の性質$(A+B)^T=A^T+B^T,(AB)^T=B^TA^T$が成り立つか確認しましょう.\smallskip\\
\begin{tabular}{llll}
    (1)$(A+B)^T$ &
    (2)$A^T+B^T$ &
    (3)$(AB)^T$ &
    (4)$B^TA^T$ 
\end{tabular}
\begin{flushright}
    (自作)
\end{flushright}
\subsubsection*{解答}
\paragraph{問題3.1}
\begin{enumerate}[label=(\arabic*)]
    \item 
    \[3\begin{bmatrix}
        3 & 2 \\ 5 & 1
    \end{bmatrix}-\begin{bmatrix}
        1 & 2 \\ 3 & -1
    \end{bmatrix}=
    \begin{bmatrix}
        9 & 6 \\ 15 & 3
    \end{bmatrix}-\begin{bmatrix}
        1 & 2 \\ 3 & -1
    \end{bmatrix}=\begin{bmatrix}
        8 & 4 \\ 12 & 4
    \end{bmatrix}\qquad\left(=4\begin{bmatrix}
        2 & 1 \\ 3 & 1
    \end{bmatrix}\right)\]
    \item $2\times1$と$2\times2$のサイズ同士の計算である.サイズが異なる行列同士の加減算はできないので,
    この計算は「計算不可能」となる.
    \item \[\begin{bmatrix}
        1 & 0 & -4
    \end{bmatrix}\]
\end{enumerate}
\paragraph{問題3.2}
\begin{enumerate}[label=(\arabic*)]
    \item \begin{align*}
        \begin{bmatrix}
            1 \\ 4
        \end{bmatrix}
        \begin{bmatrix}
            1 & 0 & 0
        \end{bmatrix}+
        \begin{bmatrix}
            2 \\ 5
        \end{bmatrix}
        \begin{bmatrix}
            0 & 1 & 0
        \end{bmatrix}+
        \begin{bmatrix}
            3 \\ 6
        \end{bmatrix}
        \begin{bmatrix}
            0 & 0 & 1
        \end{bmatrix} &=
        \begin{bmatrix}
            1 & 0 & 0 \\ 4 & 0 & 0
        \end{bmatrix}+
        \begin{bmatrix}
            0 & 2 & 0 \\ 0 & 5 & 0
        \end{bmatrix}+
        \begin{bmatrix}
            0 & 0 & 3 \\ 0 & 0 & 6
        \end{bmatrix} \\
        &=
        \begin{bmatrix}
        1 & 2 & 3 \\ 4 & 5 & 6
        \end{bmatrix}.
    \end{align*}
    掛けている行列が単位行列であることがポイントである.
    これを通して「\textbf{単位行列は要素を変えない}」ことを覚えてほしい.これが単位行列の
    最重要な性質だと考える.
    \item $3\times3$と$2\times3$のサイズ同士の計算である.定義上計算できないため「計算不可能」となる.
    \item \[R(\theta)=\begin{bmatrix}
        \cos\theta & -\sin\theta \\ \sin\theta & \cos\theta
    \end{bmatrix},\qquad R(\phi)=\begin{bmatrix}
        \cos\phi & -\sin\phi \\ \sin\phi & \cos\phi
    \end{bmatrix}\]
    \begin{align*}
        R(\theta)R(\phi) &=
        \begin{bmatrix}\cos\theta \\ \sin\theta \end{bmatrix}
        \begin{bmatrix} \cos\phi & -\sin\phi \end{bmatrix} + 
        \begin{bmatrix}-\sin\theta \\ \cos\theta \end{bmatrix}
        \begin{bmatrix} \sin\phi & \cos\phi \end{bmatrix}\\
        & = 
        \begin{bmatrix}
            \cos\theta\cos\phi & -\cos\theta\sin\phi \\ \sin\theta\cos\phi & -\sin\theta\sin\phi
        \end{bmatrix}+
        \begin{bmatrix}
            -\sin\theta\sin\phi & -\sin\theta\cos\phi \\ \cos\theta\sin\phi & \cos\theta\cos\phi
        \end{bmatrix}\\
        & =
        \begin{bmatrix}
            \cos\theta\cos\phi -\sin\theta\sin\phi & -\cos\theta\sin\phi -\sin\theta\cos\phi \\ \sin\theta\cos\phi + \cos\theta\sin\phi & -\sin\theta\sin\phi + \cos\theta\cos\phi
        \end{bmatrix}.
    \end{align*}
    もちろんこれで正解で良いが,私が出題したのには理由がある.加法定理より,
    \begin{align*}
        R(\theta)R(\phi) &=
        \begin{bmatrix}
            \cos\theta\cos\phi -\sin\theta\sin\phi & -\cos\theta\sin\phi -\sin\theta\cos\phi \\ \sin\theta\cos\phi + \cos\theta\sin\phi & -\sin\theta\sin\phi + \cos\theta\cos\phi
        \end{bmatrix}\\
        &=
        \begin{bmatrix}
            \cos(\theta+\phi) & -\sin(\theta+\phi) \\ \sin(\theta+\phi) & \cos(\theta+\phi)
        \end{bmatrix} 
    \end{align*}
    とまとめることができる.面白いことに,回転行列の積は結果として角度を足すこととなる.まるでド・モアブルの定理のようだ.
\end{enumerate}
\paragraph{問題3.3}
\begin{enumerate}[label=(\arabic*)]
    \item \[
        A+B = \begin{bmatrix}
            5 & 2 \\ 1 & 6
        \end{bmatrix}\quad \longrightarrow \quad
        (A+B)^T = \begin{bmatrix}
            5 & 1 \\ 2 & 6
        \end{bmatrix}.\]
    \item \[
        A^T + B^T =
        \begin{bmatrix}
            1 & 2 \\ 2 & 4
        \end{bmatrix}+
        \begin{bmatrix}
            4 & -1 \\ 0 & 2
        \end{bmatrix}=
        \begin{bmatrix}
            5 & 1 \\ 2 & 6
        \end{bmatrix}.
    \]
    \item \begin{align*}
        AB =
        \begin{bmatrix}
            4 & 0 \\ 8 & 0
        \end{bmatrix}+
        \begin{bmatrix}
            -2 & 4 \\ -4 & 8
        \end{bmatrix}=
        \begin{bmatrix}
            2 & 4 \\ 4 & 8
        \end{bmatrix}
        \quad \longrightarrow \quad
        (AB)^T =
        \begin{bmatrix}
            2 & 4 \\ 4 & 8
        \end{bmatrix}.
    \end{align*}
    \item \[
    B^TA^T =  \begin{bmatrix}
            4 & -1 \\ 0 & 2
        \end{bmatrix}\begin{bmatrix}
            1 & 2 \\ 2 & 4
        \end{bmatrix}=\begin{bmatrix}
            2 & 4 \\ 4 & 8
        \end{bmatrix}. \]
\end{enumerate}
\newpage